\chapter{Anomalies and global symmetries}\label{chap:anomalies}

Not every assignment of charges defines a consistent gauge theory. A chiral theory can suffer from a
\emph{gauge anomaly}: a symmetry of the classical theory that is necessarily broken by quantization. For a gauge
symmetry this is fatal. On the lattice the anomaly shows up in a simple and concrete way: the electric field
does not commute with itself.

\section{The anomaly on the lattice}
The Gauss law says that the gauge transformations are generated by $G(x) = \partial_xE - g\rho$. For a consistent
theory these generators must commute; equivalently, the electric field determined by the Gauss law must satisfy
$\comm{E(x)}{E(y)} = 0$. Using $\comm{\varphi_j}{\vartheta_k} = i\,\mathbb{1}(j\leq k)$, which follows from
$\vartheta_k = \sum_{l\leq k}P_l$, one finds for one gauge field
\begin{equation}\label{eq:EE}
  \comm{E_j}{E_k} = i\,\frac{g^2}{\pi}\,(Q^V\!\cdot Q^A)\,\operatorname{sign}(k-j),\qquad
  Q^V\!\cdot Q^A = \frac14\Bigl(\sum Q_L^2 - \sum Q_R^2\Bigr).
\end{equation}
The electric field commutes with itself if and only if
\begin{equation}\label{eq:anomaly-free}
  \sum_i Q_{L,i}^2 = \sum_j Q_{R,j}^2 ,
\end{equation}
the anomaly-cancellation condition of two-dimensional chiral gauge theories~\cite{Bertlmann1996}. The coefficient
$\sum Q_L^2 - \sum Q_R^2$ is the anomaly: it counts the mismatch of the charged left- and right-movers.
\Cref{fig:anomalous_commutator} shows \eqref{eq:EE} for several anomalous charge assignments.

\latsimfigure{anomalous_commutator}

\begin{example}
The Schwinger model ($1^2 = 1^2$) and its axial version are anomaly free. The 3-4-5-0 model has
$3^2 + 4^2 = 25 = 5^2 + 0^2$: a genuinely chiral but consistent theory. Charges $Q_L = (2,1)$, $Q_R = (0,1)$
give $5 - 1 = 4 \neq 0$ and are rejected by \code{ChiralTheory} unless \code{allow_anomalous=True}.
\end{example}

\paragraph{Several gauge fields.} For $U(1)^k$ every pair of gauge fields must be anomaly free,
\begin{equation}
  \mathcal{A}_{ab} = Q_{L,a}\cdot Q_{L,b} - Q_{R,a}\cdot Q_{R,b} = 0,
\end{equation}
including the \emph{mixed} anomalies $a\neq b$ (\code{ChiralTheory.anomaly_matrix}). In vector-axial form
$\mathcal{A}_{ab} = 2(Q^V_a\cdot Q^A_b + Q^A_a\cdot Q^V_b)$.

\paragraph{A lattice subtlety for $k>1$.} On the lattice
$\comm{E_{a,j}}{E_{b,k}} = i\frac{g_ag_b}{\pi}\bigl(x_{ab}\mathbb{1}(j\leq k) - x_{ba}\mathbb{1}(k\leq j)\bigr)$
with $x_{ab} = Q^V_a\cdot Q^A_b$. Anomaly freedom only gives $x_{ab} + x_{ba} = 0$, so a constant commutator
$i x_{ab}$ between \emph{different} gauge fields can remain. Since
$x_{ab} = (Q_{R,a}\cdot Q_{L,b} - Q_{L,a}\cdot Q_{R,b})/4$ depends on how left and right movers are paired, it is
an artefact of the regularization, absent for pairings with $Q_LQ_R^T$ symmetric.
\code{build_gauge_model} searches for such a pairing automatically.

\section{Which global symmetries survive}
Each fermion number is a global $U(1)$ symmetry of the free theory: $n_L + n_R$ of them, with charges
$t = (t_L, t_R)$. After gauging, a global symmetry survives only if it has no \emph{mixed} anomaly with the
gauge field,
\begin{equation}\label{eq:mixed}
  t_L\cdot Q_{L,a} = t_R\cdot Q_{R,a}\qquad\text{for all } a .
\end{equation}
The others are broken by the anomaly, just as the axial charge of the Schwinger model is not conserved. For an
anomaly-free theory the gauge charges themselves satisfy \eqref{eq:mixed}. The number of surviving symmetries is
$n_L + n_R - \operatorname{rank}[Q_L, -Q_R]$ (\code{ChiralTheory.global_symmetries}): one for the Schwinger model
(the gauge charge itself), three for the 3-4-5-0 model.

\paragraph{The lattice finds them by itself.} On the ring the conserved charges are the kernel directions of
$MJ$ (\cref{sec:zero-modes}). In terms of the fermion numbers, a conserved combination $t_n\cdot n + t_m\cdot m$
corresponds to $t_L\cdot N_L + t_R\cdot N_R$ with $t_n = (t_L+t_R)/2$, $t_m = (t_L-t_R)/2$, and \eqref{eq:mixed}
becomes $Q^A_a\cdot t_n + Q^V_a\cdot t_m = 0$. The tests verify that the conserved charges found numerically from
the lattice Hamiltonian span exactly the space of \eqref{eq:mixed}: the anomaly structure is not put in by hand, it
emerges from $M$.
